\section{Discussion and Conclusion}
\label{sec:conclusion}

We asked whether physical inductive biases make latent world models obey physics, and answered with a controlled anatomy: across five injection mechanisms, two training regimes, three synthetic domains, and two photorealistic-simulation benchmarks, physical structure consistently harms OOD and long-horizon prediction---even the exact gravitational form, even with perfect labels, and \emph{more} when co-trained from scratch. The failure has a mechanism, the load-bearing problem: information that is present in a shared latent is not thereby used by prediction, and structure grafted onto dimensions that carry no gradient weight decorates rather than constrains. The mechanism's falsifiable prediction held---reweighting removes the harm---but restores only parity. What moves physical robustness lies in training and data: free rollout ($2.2$--$3.6\times$, seed-robust, on synthetic and photorealistic scenes alike), horizons matched to dynamics complexity (up to $19\times$ on long-horizon photorealistic video when geometric augmentation stabilizes amplitude), and augmentation matched to the domain's true nuisances---with an explicit simple-synthetic-to-photorealistic reversal boundary, where appearance jitter flips from strongest lever to catastrophic while geometric jitter keeps helping. Alongside, we distilled four evaluation traps that had made structure look helpful and augmentation look safe.

The constructive reading is architectural. Our results do not falsify physically structured world models in general; they falsify \emph{grafting} structure onto a shared latent. Extrinsic designs---where a low-dimensional physical state is the sole pathway that prediction and decoding must route through \citep{mao2024piwm}---make physics load-bearing by construction, and our analysis predicts that this, not the physical equations themselves, is the source of their reported benefits. Testing that prediction by converting a shared-latent JEPA model to an extrinsic factorization is the direct next step; conservation-projection constraints that operate without labels are a complementary route for domains where states are unobservable.

\paragraph{Limitations.}
Our anatomy runs on one compact backbone (${\sim}10$M-parameter ViT-tiny JEPA); a frozen-encoder study across seven encoders up to $749$M parameters supports the representational conclusions, but the training-dynamics conclusions should be re-verified at scale. Headline comparisons, the reweighting result, and the collision augmentation record are three-seed; the remaining anatomy cells are single-seed, consistently signed across domains. The synthetic domains are single-mechanism by design---that is what makes the factorial clean---and the hardest open cases we report honestly: velocity-OOD resists every augmentation we tried, and deformable-object scenes remain the weakest photorealistic-simulation category for all methods. Collision is judged in latent space only, as its pixel decoder is reconstruction-limited.

\paragraph{Reproducibility.}
All training configurations, evaluation protocols, per-run logs, and the exact numbers behind every table are documented; code and data splits build on public assets (LeWM, PhyWorld, Physion, Physion++). We will release the evaluation suite, including the per-horizon/per-partition audit tools and the random-encoder calibration controls.
